\begin{animation}[id="missing-number", label="Missing Number via Sum Formula"]

\compute{
  n = 8
  given = [1, 2, 3, 4, 6, 7, 8]
  expected_sum = n * (n + 1) // 2
}

\shape{arr}{Array}{size=7, data=[1, 2, 3, 4, 6, 7, 8], labels="0..6", label="Given numbers"}
\shape{info}{Array}{size=3, data=[36, 0, 0], labels="0..2", label="expected | actual | missing"}

\step
\narrate{Ta có $n = 8$ số từ 1 đến 8, nhưng thiếu một số. Tổng kỳ vọng $= \frac{8 \times 9}{2} = 36$. Duyệt mảng và tính tổng thực tế.}

\step
\highlight{arr.cell[0]}
\apply{info.cell[1]}{value=1}
\narrate{Cộng $a[0] = 1$. Tổng thực tế $= 1$.}

\step
\recolor{arr.cell[0]}{state=done}
\highlight{arr.cell[1]}
\apply{info.cell[1]}{value=3}
\narrate{Cộng $a[1] = 2$. Tổng thực tế $= 1 + 2 = 3$.}

\step
\recolor{arr.cell[1]}{state=done}
\highlight{arr.cell[2]}
\apply{info.cell[1]}{value=6}
\narrate{Cộng $a[2] = 3$. Tổng thực tế $= 3 + 3 = 6$.}

\step
\recolor{arr.cell[2]}{state=done}
\highlight{arr.cell[3]}
\apply{info.cell[1]}{value=10}
\narrate{Cộng $a[3] = 4$. Tổng thực tế $= 6 + 4 = 10$.}

\step
\recolor{arr.cell[3]}{state=done}
\highlight{arr.cell[4]}
\apply{info.cell[1]}{value=16}
\narrate{Cộng $a[4] = 6$. Tổng thực tế $= 10 + 6 = 16$. Chú ý ta nhảy từ 4 sang 6 -- khoảng trống chính là chỗ của số 5!}

\step
\recolor{arr.cell[4]}{state=done}
\highlight{arr.cell[5]}
\apply{info.cell[1]}{value=23}
\narrate{Cộng $a[5] = 7$. Tổng thực tế $= 16 + 7 = 23$.}

\step
\recolor{arr.cell[5]}{state=done}
\highlight{arr.cell[6]}
\apply{info.cell[1]}{value=31}
\narrate{Cộng $a[6] = 8$. Tổng thực tế $= 23 + 8 = 31$.}

\step
\recolor{arr.cell[6]}{state=done}
\apply{info.cell[2]}{value=5}
\recolor{info.cell[2]}{state=good}
\narrate{Đã duyệt hết tất cả phần tử. Số bị thiếu $= 36 - 31 = 5$. Vậy số bị thiếu là $\mathbf{5}$.}
\end{animation}
